\section*{الفصل \ref{ch:b1:gas-exchange} --- التبادل الغازي التنفسي}

\begin{solution}{exo:b1:gas-exchange:1}
$\dot V = K(A/x)(P_1 - P_2)$: مساحة أكبر، وحاجز أرق، \omterm{def:b1:gas-exchange:partial}{وضغط جزئي}
خارجي أعلى (التهوية)، \omterm{def:b1:gas-exchange:partial}{وضغط جزئي} داخلي أخفض (الترواء).
\end{solution}

\begin{solution}{exo:b1:gas-exchange:2}
أكسجين أقل في اللتر بثلاثين مرة (\qty{7}{} في مقابل
\qty{210}{mL/L})؛ وأكثف بثماني مئة مرة وألزج بخمسين مرة، فضخه
باهظ؛ والأكسجين ينتشر فيه أبطأ عشرة آلاف مرة، فوجب تحريك الوسط
عبر السطح كله.
\end{solution}

\begin{solution}{exo:b1:gas-exchange:3}
بتيار موافق: يغادر الماء والدم كلاهما عند \qty{12}{kPa}؛
والاستخلاص $(21 - 12)/21 = \qty{43}{\%}$. وبتيار معاكس: يغادر
الماء عند \qty{4}{kPa} والدم عند \qty{18}{kPa}؛ والاستخلاص
$(21 - 4)/21 =
\qty{81}{\%}$.
\end{solution}

\begin{solution}{exo:b1:gas-exchange:4}
منحلًا (\qty{7}{\%})، ومرتبطًا بالزمر الأمينية في الهيموغلوبين
(\qty{23}{\%})، وبيكربونات في البلازما (\qty{70}{\%}).
\end{solution}

\begin{solution}{exo:b1:gas-exchange:5}
الدقيقية \qty{6.75}{L/min}؛ والسنخية $15\times 0.3 = \qty{4.5}{L/min}$؛
والأخذ $4.5\times 0.07 = \qty{0.32}{L/min}$. وبالأنفاس السطحية:
الدقيقية \qty{6.75}{L/min} كما هي، لكن السنخية $30\times 0.075 =
\qty{2.25}{L/min}$ والأخذ \qty{0.16}{L/min} --- أي النصف، مقابل
عمل التنفس نفسه؛ فالحيّز الميت يُدفع ثمنه في كل نفس.
\end{solution}

\begin{solution}{exo:b1:gas-exchange:6}
$2/(6\times 0.75) = \qty{0.44}{L}$ في الساعة، أي اثنان وعشرون مثلًا
من حجمها البالغ \qty{20}{mL}.
\end{solution}

\begin{solution}{exo:b1:gas-exchange:7}
تهبط $A/x$ بالمعامل $2\times 4 = 8$: أي \qty{31}{mL/min}. ولاستعادة
\qty{250}{}، يجب أن يكون الفرق \qty{64}{kPa} --- وهو مستحيل، إذ لا
يستطيع الضغط السنخي أن يتجاوز \qty{13}{kPa} في الهواء ولا يستطيع
الدم الوريدي أن يهبط دون الصفر؛ ولا يستطيع المريض رفعه إلا بتنفس
الأكسجين، وحتى عندئذ لا يرفعه ثمانية أمثال.
\end{solution}

\begin{solution}{exo:b1:gas-exchange:8}
\omterm{prop:b1:gas-exchange:transport}{أنهيدراز الكربونيك} داخل \omterm{def:b1:cell-unit-of-life:cell}{الخلية}، فتُصنع البيكربونات هناك وتتراكم؛
ثم تنتشر خارجًا مع تدرجها عبر مبادل يُدخل كلوريدًا مقابل كل
بيكربونات، فيبقى الشحن متوازنًا. ولولا هذه المبادلة لبقيت
البيكربونات في الداخل مع شريك بروتونها، فرفعت أسمولية \omterm{def:b1:cell-unit-of-life:cell}{الخلية}،
ودخلها الماء: فتنتفخ \omterm{def:b1:cell-unit-of-life:cell}{الخلية} في \omterm{def:b1:body-plans-tissues:tissue}{الأنسجة} (وهي تنتفخ قليلًا كما هي).
\end{solution}

\begin{solution}{exo:b1:gas-exchange:9}
الرغبة في التنفس تأتي من ارتفاع $\mathrm{CO_2}$. وفرط التهوية يخفض
$\mathrm{CO_2}$ إلى ما دون الطبيعي بكثير دون أن يضيف أكسجينًا يذكر
(فالهيموغلوبين كان مشبعًا أصلًا)؛ وفي أثناء الغوصة يهبط الأكسجين
إلى مستوى فقدان الوعي قبل أن يعود $\mathrm{CO_2}$ صاعدًا إلى
المستوى الذي يفرض نفسًا. وبلا فرط تهوية يبلغ $\mathrm{CO_2}$ ذلك
المستوى والأكسجين ما يزال وافرًا، فيصعد الغائص.
\end{solution}

\begin{solution}{exo:b1:gas-exchange:10}
$A = \pi(10^{-5})^2 = \qty{3.1e-10}{m^2}$؛ و$\Delta c = JL/DA =
10^{-9}\times 10^{-3}/(2\times 10^{-5}\times 3.1\times 10^{-10}) =
\qty{160}{mol/m^3}$ --- أي أكثر بكثير من \qty{8.7}{} المتاحة، فأنبوب
بهذا الضيق لا يستطيع إمداد العضلة بالانتشار عند \qty{1}{mm}؛
ويلزم أنبوب أوسع (\qty{100}{\micro m}: $\Delta c = \qty{6.4}{mol/m^3}$، وهو
ممكن بالكاد) أو تهوية فاعلة. وعند \qty{10}{mm} يصير المطلوب عشرة
أمثال ذلك. ولما كان طول الأنبوب ينمو مع حجم الجسم والطلب ينمو مع
الحجم، حدّ الانتشار عبر \omterm{prop:b1:gas-exchange:designs}{القصبات الهوائية} حجم الحشرات ببضعة
سنتيمترات، ولا يزيد إلا بالضخ.
\end{solution}

\begin{solution}{exo:b1:gas-exchange:11}
الأسماك: الماء في اتجاه واحد والدم معاكس له (\omterm{def:b1:gas-exchange:gill}{تيار معاكس})،
والاستخلاص \qty{80}{\%}. والطيور: الهواء في اتجاه واحد عبر
\omterm{def:b1:flowering-plant-organization:tissues}{القصيبات} الموازية، والدم عمودي عليه (تيار متقاطع)، والاستخلاص
وسط لكنه أعلى من استخلاص الثدييات، ولا يتلقى \omterm{prop:b1:organism-environment:surfaces}{سطح التبادل} هواءً
بائتًا أبدًا. أما \omterm{def:b1:gas-exchange:lung}{الرئة} المدّية فتخلط الهواء العذب بالغاز
المتبقي، فلا يرى \omterm{prop:b1:organism-environment:surfaces}{سطح التبادل} في أحسن الأحوال إلا الخليط السنخي،
ولا يستطيع الدم أن يتوازن إلا معه: فمهما يكن السطح لا يستطيع الدم
الشرياني أن يتجاوز الضغط السنخي، ولا يستطيع الضغط السنخي أن يقارب
الضغط الملهَم ما بقي حجم متبقٍّ.
\end{solution}

\begin{solution}{exo:b1:gas-exchange:12}
$\mathrm{CO_2}$ هو الإشارة الأفضل لأن مستواه الشرياني يستجيب
للتهوية استجابة مباشرة وحادة (انصف التهوية فيتضاعف)، ولأنه يُقاس
بدقة عبر pH بالمستقبلات المركزية، ولأن الأكسجين، وهو جالس على قمة
منحني الهيموغلوبين المسطحة، يتغير قليلًا في المدى الطبيعي
للتهوية: فمنظّم حراري يقوم على $\mathrm{CO_2}$ يبقي الأكسجين
صحيحًا كناتج ثانوي. ويخفق الترتيب حين يهبط الأكسجين مستقلًا عن
$\mathrm{CO_2}$ --- أي في المرتفعات، حيث يكون الأكسجين الملهَم
منخفضًا وإنتاج $\mathrm{CO_2}$ غير منخفض، فتقول إشارة
$\mathrm{CO_2}$ «تنفس أقل» بينما يقول الأكسجين «تنفس أكثر»؛ ولا
تتجاوزها إلا الأجسام السباتية حين يشتد نقص الأكسجين، وتستغرق
التسوية (فرط تهوية مع قلاء) أيامًا من الضبط الكلوي حتى تستقر.
\end{solution}

\begin{solution}{pb:b1:gas-exchange:1}
\textbf{1.} الهواء $210/22.4 = \qty{9.4}{mmol/L}$؛ والماء $7/22.4 =
\qty{0.31}{mmol/L}$.
\textbf{2.} 30.
\textbf{3.} الهواء: $1.2/9.4 = \qty{0.13}{g}$؛ والماء: $1000/0.31 =
\qty{3200}{g}$ --- أي كتلة أكبر بمقدار \num{25000} مرة.
\textbf{4.} $210/7 = \qty{30}{L}$ من الماء.
\textbf{5.} الماء أثقل وأفقر بالأكسجين من أن يُدفع داخلًا وخارجًا؛
بل يجب تمريره مرة واحدة، وهو ما يتيح أيضًا \omterm{def:b1:gas-exchange:gill}{التيار المعاكس}؛ أما
الهواء فرخيص بما يكفي لتحريكه مرتين.
\textbf{6.} الدقيقية $12\times 0.5 = \qty{6}{L/min}$؛ والسنخية
$12\times 0.35 = \qty{4.2}{L/min}$.
\textbf{7.} $150/500 = \qty{30}{\%}$ من كل نفس، ومن التهوية الدقيقية.
\textbf{8.} الموصل $4.2\times 0.21 = \qty{0.88}{L/min}$؛ والمأخوذ
$4.2\times 0.07 = \qty{0.29}{L/min}$، وهو قريب من \qty{250}{mL/min}.
\textbf{9.} $0.25/(6\times 0.21) = \qty{20}{\%}$ (أي \qty{28}{\%} مما
يبلغ \omterm{def:b1:gas-exchange:lung}{الأسناخ}).
\textbf{10.} $0.14\times(101 - 6.3) = \qty{13.3}{kPa}$.
\textbf{11.} $250/5 = \qty{50}{mL/L}$؛ والوريدي \qty{150}{mL/L}،
مشبع \qty{75}{\%}.
\textbf{12.} $3000/20 = \qty{150}{mL/L}$؛ والوريدي \qty{50}{mL/L}،
مشبع \qty{25}{\%}.
\textbf{13.} $50/60 = \qty{0.83}{mL/min}$.
\textbf{14.} $0.83/(7\times 0.8) = \qty{0.149}{L/min}$، أي \qty{8.9}{L}
في الساعة.
\textbf{15.} $0.83/(7\times 0.43) = \qty{0.28}{L/min}$: \omterm{def:b1:gas-exchange:gill}{فالتيار
المعاكس} ينصّف الماء الواجب ضخه، وينصّف العمل.
\textbf{16.} \qty{8.9}{kg} من الماء في الساعة؛ والإنسان يحرك $6\times
60\times 1.2 = \qty{430}{g}$ من الهواء في الساعة --- أي جزء من عشرين
من الكتلة، مقابل ثلاث مئة مثل من الأكسجين ($15\,000$ في مقابل
\qty{50}{mL} في الساعة).
\textbf{17.} $100\times(0.95 - 0.15) = \qty{80}{mL/L}$؛ والناتج القلبي
$0.83/0.080 = \qty{10.4}{mL/min}$.
\textbf{18.} الماء إلى الدم $149/10.4 = 14$؛ والهواء إلى الدم $4.2/5 = 0.84$.
فعلى السمكة أن تمرر أربعة عشر حجمًا من الماء لكل حجم من الدم لأن
كلًّا منهما يحمل من الأكسجين قدرًا ضئيلًا؛ أما الإنسان فنحو واحد.
\textbf{19.} الأخذ $\times 1.5$ من ماء يحمل $5.5/7 = 0.79$ من
المقدار: فالتهوية $\times 1.9$؛ وكلفة التنفس، وهي متناسبة تقريبًا
مع الدفق (أو أكثر، إذ ترتفع أسرع من الخطية)، تنتقل من \qty{10}{\%}
إلى \qty{13}{\%} على الأقل من أخذ أكبر --- فالسمكة تكدّ أكثر في
التنفس في الماء الدافئ.
\textbf{20.} عند طرف دخول الماء: $21 - 18 = \qty{3}{kPa}$؛ وعند طرف
خروجه: $4 - 3 = \qty{1}{kPa}$: أي موجب عند الطرفين.
\textbf{21.} عند طرف الدخول $21 - 3 = \qty{18}{kPa}$؛ وعند طرف الخروج $12 - 12 =
0$: فقد تعادل المائعان ولم يبق تدرج، مع أن الماء ما يزال يحمل
\qty{57}{\%} من أكسجينه.
\textbf{22.} تشبه التيار الموافق: فالدم يتوازن مع غاز سنخي واحد
جيد الخلط ويغادر عند ضغطه؛ ولا يستطيع تجاوز \qty{13.3}{kPa} لأنه
لا غاز أعذب في المجرى يلقاه، كما هي الحال في \omterm{def:b1:gas-exchange:gill}{التيار المعاكس}.
\textbf{23.} السمكة $2000/(2\times 10^{-4}) = \num{1e7}$ (cm$^2$ لكل cm)؛
والإنسان $10^6/(5\times 10^{-5}) = \num{2e10}$: أي نسبة \num{2000}. والأخذان
$250/0.83 = 300$. فالسمكة تأخذ من الأكسجين سبعة أمثال ما يأخذه
الإنسان لوحدة $A/x$: فلا بد أن يكون فرق الضغط الوسطي عبر صفيحتها
أكبر --- وهو ما يوفره \omterm{def:b1:gas-exchange:gill}{التيار المعاكس}، إذ يحفظ تدرجًا على السطح كله
حيث يكون تدرج \omterm{def:b1:gas-exchange:lung}{الرئة} المدّية صغيرًا قرب نهاية التعادل.
\textbf{24.} $\mathrm{CO_2}$ أذوب في الماء من الأكسجين بخمس وعشرين
مرة: فالماء يحمله بعيدًا بسرعة إنتاجه، ويبقى $\mathrm{CO_2}$ في دم
السمكة منخفضًا جدًا.
\textbf{25.} \omterm{def:b1:gas-exchange:gill}{الغلصمة} نحو \qty{80}{\%}، \omterm{def:b1:gas-exchange:lung}{والرئة} نحو \qty{20}{\%} من
الأكسجين في الوسط الملهَم؛ والجريان في اتجاه واحد وترتيب الماء
والدم \omterm{def:b1:gas-exchange:gill}{بتيار معاكس} هما ما يصنع الفرق.
\end{solution}
